% =====================================================================
%  学習指導案：情報科（情報I）
%  ---------------------------------------------------------------------
%  単元：アクセシブルなデザイン（小単元）
%  本時：アクセシブルなデザインの実践（ペルソナワーク）
%  形式：オンライン非同期型授業
% =====================================================================
\documentclass[uplatex,a4paper,11pt]{jsarticle}

\usepackage{geometry}
\geometry{top=20mm,bottom=20mm,left=20mm,right=18mm}

\usepackage{array}
\usepackage{longtable}
\usepackage{enumitem}

\setlength{\tabcolsep}{4pt}
\renewcommand{\arraystretch}{1.35}

\newcommand{\itemhead}[2]{\par\medskip\noindent\textbf{#1.\quad #2}\par\nopagebreak\smallskip}

\pagestyle{plain}

\begin{document}

\begin{center}
  {\LARGE \textbf{情報科 学習指導案}}
\end{center}

\begin{flushright}
  \begin{tabular}{ll}
    指導教諭   & 〇〇　〇〇　　印 \\
    教育実習生 & 吉岡　優　　印 \\
  \end{tabular}
\end{flushright}

\bigskip

\noindent
1.\quad 日\hspace{1zw}時：令和8年5月26日（火）　第9,10校時相当（オンライン非同期型授業／想定学習時間50分） \par
\noindent
2.\quad 実施学級：東京都立　昭和高等学校　第1学年H組　在籍生徒数　38名 \par
\noindent
3.\quad 単\hspace{1zw}元：『最新情報Ⅰ』（実教出版）　第〇章　コミュニケーションと情報デザイン　第〇節　アクセシブルなデザイン \par


\itemhead{4}{単元観}
\hspace{1zw}本単元は、高等学校学習指導要領（平成30年告示）第2章第10節「情報」第1款情報Ⅰ「(2)コミュニケーションと情報デザイン」を受けて設定したものであり、情報デザインが人や社会に果たしている役割を理解し、目的や状況に応じて受け手に分かりやすく情報を伝える表現方法を身に付けることを目指している。\par
\hspace{1zw}現代社会では、Webサイト・アプリ・公共サイン等あらゆる情報が画面越しに届けられるようになり、その情報が誰に届き、誰に届かないかという問題が顕在化している。WebAIM Million Report（2024）によれば、世界のトップ100万サイトの95\%以上に何らかのアクセシビリティ上の問題が存在する。このような状況下で、情報デザインを単なる「見た目を整える技術」として習得させるだけでは不十分である。\par
\hspace{1zw}本単元では、色覚多様性・言語・状況など多様な利用者の存在を意識し、「誰のためにデザインするのか」という根源的な問いに生徒を向き合わせる。情報技術を扱う者としての社会的責任を理解し、他者への想像力を育むことを通して、情報社会における主体的な担い手としての姿勢を育てたい。\par

\itemhead{5}{学級・生徒観}
\hspace{1zw}第1学年H組は38名（男子18名・女子20名）。情報Ⅰの前単元で情報デザインの基礎（抽象化・可視化・構造化）を学習済みであり、HTML/CSSの簡単な記述経験も有する。東京都の政策により生徒には一人一台Surfaceが配付されており、全員が自宅から学校LMSにアクセス可能な環境にある。\par
\hspace{1zw}学習姿勢には注意をおくべきである。授業中スマホを机の下で触っていることや、寝ている生徒がちらほらみられる。しかし、グループでの話し合いには積極的に取り組む生徒が多いが、自分の意見を文章で深く書くことに苦手意識を持つ生徒が一定数いる。色覚多様性については、男子生徒のうち統計的に1〜2名程度が該当する可能性があり、配色を扱う本時では当該生徒への配慮と、他生徒の気づきの場としての両側面からの設計が必要である。\par
\hspace{1zw}非同期型授業の実施経験は本年度3回目であり、自学管理には慣れているが、演習が長くなると集中力が低下する傾向が見られる。AI（生成AI）の利用経験は約7割の生徒にあり、AIアシスタント機能の操作には抵抗が少ないと予想される。\par

\itemhead{6}{指導観・教材観}
\hspace{1zw}本時は小単元「アクセシブルなデザイン」（5時間扱い）の第4時に位置付ける。第1〜3時で情報デザインの基本・色彩理論・WCAGコントラスト比の知識を習得した上で、本時はそれらを統合的に実践する場として設計する。\par
\hspace{1zw}本単元はVygotskyの最近接発達領域（ZPD）理論に基づき、段階的なスキャフォールディング（足場かけ）を採用している。本時（第4時）では教師が用意した共通ペルソナ「祖父」を全生徒で扱い、多様性への気づきと配色設計の手順を確実に習得させる。次時（第5時）では生徒自身が身近な誰かを想定したペルソナを設定し、本時で学んだ知識・技能を自立的に応用する展開とする。これにより、初学者が安心して取り組める基盤を保ちながら、最終的には学習者の主体性を引き出す構造を実現する。\par
\hspace{1zw}共通ペルソナ「祖父」は、以下のプロフィールで設定する。「あなたの祖父（68歳）。色覚多様性（赤緑色弱）があり、老眼も進んでいる。スマートフォンは持っているが操作はゆっくり。孫であるあなたの写真を見るのが何より楽しみ。」色覚シミュレータ・WCAGコントラスト比エディタという教材構造と完全に整合し、感情的なつながり（家族）を通じて生徒の動機を引き出せる題材である。\par
\hspace{1zw}なお、学級・生徒観で示した「文章で深く書くことに苦手意識を持つ生徒」「演習が長くなると集中力が低下する傾向」に対応するため、本時では(1)動画を各5分以内に区切ること、(2)演習を細かいステップに分けること、(3)振り返り記述には書き出しのヒントを提示すること、の3点で配慮する。\par
\hspace{1zw}なお、本時で扱う「祖父」は身体的特性に基づくアクセシビリティであるが、Microsoftが提唱するInclusive DesignにおけるPersona Spectrum（Permanent／Cultural／Situational）の3軸が存在することは動画②内で紹介し、次時（第5時）の自由ペルソナ設定の参考とする。これにより「アクセシビリティ＝障害者対応」という狭い理解から、「すべての人が状況によって使いづらさを経験する」という広い理解へと拡張させる。\par
\hspace{1zw}教材は反転学習の構造を採用し、動画によるインプットとWebアプリによる演習・ピアラーニングを組み合わせる。AIアシスタントによる壁打ち機能と、作品ライブラリによる蓄積型ピア閲覧を通じて、非同期型授業における対話的学びを実現する。作品ライブラリには教師が用意したサンプル作品と生徒の提出作品が同一の場に蓄積され、本時（第4時）ではサンプル作品を、次時（第5時）の相互評価では全作品を活用する。教材内で扱う全ての素材は自作またはライセンスフリーであり、情報Ⅰの「著作権」単元の模範教材としても機能するよう設計した。\par

\itemhead{7}{単元の目標}
\begin{enumerate}[label=(\arabic*),leftmargin=2.8em,itemsep=3pt,topsep=2pt]
  \item 情報デザインの基本原則およびアクセシビリティの考え方（ユニバーサルデザイン、WCAG等）を理解し、配色のコントラスト比を算出・評価することができる。（知識・技能）
  \item 多様なユーザーの特性（身体的・文化的・状況的）を踏まえ、WCAG等の基準と他者への想像力を統合してデザインの選択を多面的に説明することができる。（思考・判断・表現）
  \item 多様な利用者の存在を意識し、他者への想像力をもって情報デザインに取り組もうとする。（主体的に学習に取り組む態度）
\end{enumerate}

\itemhead{8}{単元の指導計画（5時間扱い）}
\nopagebreak
\hspace{1zw}※本単元は大単元「コミュニケーションと情報デザイン」（15時間扱い）の終盤に位置付く小単元「アクセシブルなデザイン」（5時間扱い）として構成する。\par
\smallskip
{\renewcommand{\arraystretch}{1.4}
\begin{longtable}{|>{\centering\arraybackslash}m{16mm}|p{48mm}|p{96mm}|}
  \hline
  \multicolumn{1}{|c|}{時} & \multicolumn{1}{c|}{学習活動} &
  \multicolumn{1}{c|}{評価規準【観点】（評価方法）} \\ \hline
  \endfirsthead
  \hline
  \multicolumn{1}{|c|}{時} & \multicolumn{1}{c|}{学習活動} &
  \multicolumn{1}{c|}{評価規準【観点】（評価方法）} \\ \hline
  \endhead
  1 &
    情報デザインの基本とユニバーサルデザインの考え方\newline pp.〇〇--〇〇 &
    ・情報デザインの基本原則とユニバーサルデザインの理念を説明している。\newline 【知識・技能】（ワークシート） \\ \hline
  2 &
    色彩理論と視覚効果\newline pp.〇〇--〇〇 &
    ・色相・明度・彩度の基本および色彩の心理的効果を理解している。\newline 【知識・技能】（小テスト） \\ \hline
  3 &
    アクセシビリティの基礎（色覚多様性とWCAG）\newline pp.〇〇--〇〇 &
    ・色覚多様性の事実とWCAGコントラスト比の意味を説明している。\newline 【知識・技能】（ワークシート） \\ \hline
  4\newline（本時） &
    アクセシブルなデザインの実践（共通ペルソナ「祖父」による配色設計）\newline pp.〇〇--〇〇 &
    ・WCAGに基づき共通ペルソナの特性に応じた配色を設計している。\newline 【知識・技能・思考・判断・表現】（成果物・振り返り記述） \\ \hline
  5 &
    自由ペルソナによる応用と作品の相互評価\newline pp.〇〇--〇〇 &
    ・自分で設定したペルソナに応じた配色を設計し、他者の作品から多様な視点を学んでいる。\newline 【主体的に学習に取り組む態度】（成果物・コメント・レポート） \\ \hline
\end{longtable}}

\itemhead{9}{本時の目標}
\begin{itemize}[label={○},leftmargin=2.2em,itemsep=3pt,topsep=2pt]
  \item WCAGコントラスト比を用いて配色のアクセシビリティを評価し、共通ペルソナ「祖父」の特性（色覚多様性・老眼）に応じた配色を設計することができる。【知識・技能】
  \item 自分以外の他者の視点に立ち、デザインの選択理由を多面的に説明することができる。【思考・判断・表現】
\end{itemize}

\itemhead{10}{本時の授業展開}
\nopagebreak
{\renewcommand{\arraystretch}{1.4}
\begin{longtable}{|>{\centering\arraybackslash}m{14mm}|p{72mm}|p{68mm}|}
  \hline
  \multicolumn{1}{|c|}{\shortstack{段階\\（時間）}} &
  \multicolumn{1}{c|}{学習活動（○指示・発問）} &
  \multicolumn{1}{c|}{指導上の留意点／☆評価の観点} \\ \hline
  \endfirsthead
  \hline
  \multicolumn{1}{|c|}{\shortstack{段階\\（時間）}} &
  \multicolumn{1}{c|}{学習活動（○指示・発問）} &
  \multicolumn{1}{c|}{指導上の留意点／☆評価の観点} \\ \hline
  \endhead
  導入\newline10分 &
    ・LMSにログインし、本時の教材ページにアクセスする。\newline
    ・イントロ動画（3分）を視聴する。\newline
    ○（基発問）「みなさんが普段スマートフォンやWebサイトを使っていて、『見にくいな』『使いにくいな』と感じた経験はありませんか？」\newline
    ・自身の経験を1〜2文で記入する。\newline
    ・教材ページに表示される橋渡し文（「もし自分の家族や友人が、自分とは違う見え方をしていたら、その人にはどう見えているでしょうか。今日は、ある架空の祖父を題材にして考えていきます。」）を読む。\newline
    ・祖父紹介動画（5分）を視聴し、孫として家族写真のアルバムサイトをデザインするという文脈を確認する。\newline
    ○（指示）「今日は、あなたの祖父にとってのアルバムサイトを考えます。祖父がどんな人なのかを知った上で、次の動画で祖父の見え方をくわしく学びましょう。」 &
    ・教師は配信開始時刻に教師管理画面でログイン状況を確認し、未着手の生徒には10分後にリマインドを送信する。\newline
    ・基発問はイントロ動画内でも投げかけ、生徒は教材ページの入力欄に経験を記入する二重構造とする。本時の思考の外枠を決める導入問いに留め、核心となる問いは主発問（展開段階）で扱う。\newline
    ・共通ペルソナ「祖父」（68歳・色覚多様性・老眼）の人物紹介は導入の祖父紹介動画で行い、見え方の詳細は展開の動画①で扱う。\newline
    ・基発問への応答後、教材ページに橋渡し文を自動表示する設計とする。\newline
    ・教師制作のイントロ動画は3分以内とし、本時の流れを簡潔に示す。\newline
    ・ペルソナは架空の人物として設定していることを動画内で明示し、実在の家族の個人情報を入力しないよう注意喚起する。\newline
    ・色覚多様性を持つ生徒には、事前に「自分の見え方」「他者の見え方」双方を体験できる選択肢を案内しておく。\newline
    ☆共通ペルソナの状況を理解しようとしている。（教師管理画面で進捗を確認） \\ \hline
  展開\newline32分 &
    ・動画①「祖父の見え方を知ろう（第3時の復習＋ケーススタディ）」（5分）を視聴する。\newline
    ・演習①：色弱シミュレータで、祖父の「見え方・困りごと」を体験し、気づきを1〜2文で記入する（5分）。\newline
    ・動画②「WCAGとデザインの原則」（5分）を視聴する。\newline
    ○（主発問）「ここまで色覚多様性とWCAGの数値を学んできました。でも、祖父は本当に『数値だけ』で見やすくなるのでしょうか。数字を満たしたデザインが、必ずしも『その人に届く』とは限りません。では、本当に祖父に届くデザインとは、どのようなものだと思いますか？」\newline
    ・演習②：祖父のアルバムサイトの配色・レイアウトを配色エディタで設計する。WCAGコントラスト比をリアルタイム確認しながら、作品を完成させて提出する（12分）。詰まった場合はAIアシスタントに相談する。\newline
    ・作品ライブラリで、祖父をテーマにしたサンプル作品3つ（機能優先・感情優先・対話派の3軸）を閲覧し、各作品に気づきコメントを1文記入する（5分）。 &
    ・動画①は第3時で学んだ色覚多様性の知識を、祖父というペルソナに当てはめて具体化する位置付けで制作する。\newline
    ・動画②内ではPersona Spectrum（Permanent／Cultural／Situational）の概念にも触れ、次時の自由ペルソナ設定の基礎とする。\newline
    ・主発問は動画②の末尾で投げかけ、教材ページにテキストで再掲して、演習②の設計意図に反映させる二重構造とする。\newline
    ・教師は教師管理画面で5分おきに各生徒の進捗を確認する。動画視聴率が30\%未満で停滞している生徒には個別チャットで声をかける。\newline
    ・演習②開始から15分経過しても提出がない生徒には、AI壁打ち機能の使用を促す通知を送る。\newline
    ・AIアシスタントは断定的な答えを避け、「ペルソナの視点ではどう見えるか」を問いかける形で支援する設計とする。\newline
    ・通信トラブル時はLMSに用意した紙ベースのフォールバック教材（PDF）に切り替えるよう案内する。\newline
    ・他者の作品にネガティブなコメントをしないことを動画内および教材内に明記する。\newline
    ☆WCAGコントラスト比を算出し作品に適用している。（成果物・自動算出値）\newline
    ☆祖父の特性を踏まえた設計意図を説明している。（成果物の設計意図メモ） \\ \hline
  整理\newline8分 &
    ・まとめ動画（2分）を視聴する。\newline
    ○（指示）「次の3つの問いに、自分の言葉で答えてください。書き出しのヒントを参考にして構いません。」\newline
    （1）祖父の見え方・困りごとを想像できましたか。\newline
    　　（ヒント：「シミュレータを体験して〜と感じた」）\newline
    （2）自分のデザインで何を最も大切にしましたか。\newline
    　　（ヒント：「私は〜を優先した。なぜなら〜」）\newline
    （3）次回は自分で設定したペルソナで設計します。誰のためにデザインしたいと思いましたか。\newline
    　　（ヒント：「次は〜さんのために設計したい。その人は〜」）\newline
    ・参考文献（WebAIM Million Report、CUDO等）へのリンクを確認する。 &
    ・教師制作のまとめ動画の最後で、次時の自由ペルソナ設定について予告する設計とする。\newline
    ・振り返りは正解を求めず、自分の言葉で記述するよう促す。\newline
    ・3つの問いは段階的に、本時の祖父→自分のデザイン→次時への動機、と思考が展開する構造になっていることを意識させる。\newline
    ・文章記述に苦手意識のある生徒のため、書き出しのヒントを各問いに付して認知負荷を下げる。\newline
    ・教師は配信終了後、全生徒の振り返り記述を読み、個別文章評価コメントを48時間以内に返信する。\newline
    ☆他者への想像力を自分の言葉で言語化している。（振り返り記述）\newline
    ☆次時の自由ペルソナ設定への意欲を示している。（記述内容） \\ \hline
\end{longtable}}

\bigskip

\noindent\textbf{【期待される振り返り記述の例】}

\hspace{1zw}「色覚多様性のシミュレータを体験してみて、普段自分が見えていると思っている世界が、祖父には全く違って見えることに驚いた。デザインを作る時、WCAGコントラスト比4.5:1以上という数値の基準を満たすことは最低限の責任だと感じた。一方で、数値を満たしただけのデザインは『無難に正しいだけ』で、祖父にとって楽しい時間を生み出すには、もう一歩想像力が必要だと思った。次回は自分の妹のためにデザインしてみたい。妹は小学生で、漢字が読めないし長文も苦手だ。今日学んだ『相手を想像する姿勢』を、別の相手で試してみたい。」

\bigskip

\noindent\textbf{【補足】指導上の特記事項}

\begin{itemize}[leftmargin=2em,itemsep=2pt,topsep=2pt]
  \item \textbf{非同期型授業特有の対応：}対面授業における机間巡視・板書・一斉発問は、それぞれ教師管理画面での進捗監視・教材ページのデザイン・動画内ナレーション及び教材内テキストに置き換えて設計している。教師は配信時間中、管理画面で全生徒の進捗を継続的に監視し、個別対応を行う。
  \item \textbf{使用するリソース：}
    \begin{itemize}[leftmargin=1em,topsep=0pt,itemsep=0pt]
      \item 動画①「祖父の見え方を知ろう（第3時の復習＋ケーススタディ）」（5分／教師制作）
      \item 動画②「WCAGとデザインの原則」（5分／教師制作）
      \item 動画③「本時のまとめ＋次時の予告」（2分／教師制作）
      \item Webアプリ：配色エディタ＋色弱シミュレータ（教師制作）
      \item 作品ライブラリ：教師作成のサンプル作品3つ（機能優先・感情優先・対話派の3軸）と生徒の提出作品が同一の場に蓄積される。本時はサンプル作品を閲覧し、次時（第5時）の相互評価では全作品を活用する。生徒作品の公開は本人及び保護者の同意を得たもののみとする。
      \item AIアシスタント：Claude Haiku 4.5を内部API利用、プロンプト設計済み
    \end{itemize}
  \item \textbf{参考文献（生徒向け）：}
    \begin{itemize}[leftmargin=1em,topsep=0pt,itemsep=0pt]
      \item WebAIM Million Report（\texttt{https://webaim.org/projects/million/}）
      \item カラーユニバーサルデザイン機構 CUDO（\texttt{https://cudo.jp/}）
      \item WCAG 2.2 日本語訳（\texttt{https://waic.jp/translations/WCAG22/}）
    \end{itemize}
  \item \textbf{個人情報保護：}ペルソナは架空人物として設定し、実在の家族・友人の情報を入力させない。提出された作品はクラス内のみ閲覧可能とし、外部公開は行わない。
  \item \textbf{著作権配慮：}教材内の全素材は自作またはライセンスフリーであり、参考文献は公的機関の一次ソースに直リンクする。引用4要件（公表性・区別性・主従関係・出典明示）を遵守し、本教材自体が情報Ⅰの著作権単元の模範となるよう設計している。
  \item \textbf{進度差への対応：}早く演習を終えた生徒には発展課題として「祖父の別の特性（例：手の動かしにくさ、聴力の低下）にも注目してデザインを見直す」を提示する。時間がかかる生徒にはテンプレートを選択可能とする。なお、自由なペルソナ設定は次時（第5時）で扱うため、本時は祖父1人での深掘りに集中させる。
  \item \textbf{学習困難への配慮：}全動画に字幕を付与。配色エディタはキーボード操作対応。色覚多様性を持つ生徒には「自分の見え方」と「他者の見え方」双方のシミュレートを可能にし、当事者性を活かした学びを提供する。
\end{itemize}

\end{document}
